\documentclass[11pt,a4paper]{scrartcl}
\usepackage[a4paper, top=20mm, bottom=20mm, left=22mm, right=22mm]{geometry}
\usepackage{lmodern}
\usepackage[T1]{fontenc}
\usepackage[utf8]{inputenc}
\usepackage[ngerman]{babel}
\usepackage{microtype}
\usepackage{parskip}
\usepackage{enumitem}
\usepackage[hidelinks]{hyperref}
\pagestyle{empty}

\begin{document}

\begin{center}
  {\Large\bfseries Kurzfassung: Erste Constrained-Benchmark-Ergebnisse}\\[4pt]
  {\normalsize 4.~März~2026}
\end{center}

\bigskip

\subsection*{Was wir gemacht haben}

100 RocksDB-Durchläufe in Docker mit 256\,MB RAM, 1 CPU, kein Swap.
Drei Parameter variiert: Schreibpuffer (4--128\,MB), Block-Cache (8--128\,MB),
Kompression (keine/snappy/lz4/zstd). Workload: gleichzeitiges Lesen und
Schreiben (\texttt{readwhilewriting}). 58 erfolgreich, 42 vom OOM-Killer
beendet.

\subsection*{Was wir schon sehen}

\begin{enumerate}[nosep]
  \item \textbf{Schreibpuffer dominiert den Speicherverbrauch.}
        Ab 64\,MB Schreibpuffer stürzt RocksDB in fast allen Fällen ab.
        Sicherer Bereich: Schreibpuffer $\leq$\,32\,MB, Cache $\leq$\,64\,MB.

  \item \textbf{Ohne Kompression geht es nicht.}
        Unkomprimiert: 2.000~ops/sec, 50\,\% Stall-Zeit.
        Mit zstd: 3.600~ops/sec, 32\,\% Stall-Zeit.
        Die I/O-Einsparung überwiegt die CPU-Kosten deutlich.

  \item \textbf{17-facher Leistungsunterschied durch Konfiguration allein.}
        Beste Konfig: 8.900~ops/sec. Schlechteste: 518~ops/sec.
        Gleiche Hardware, gleiche Daten -- nur andere Parametereinstellungen.

  \item \textbf{Write-Stalls sind massiv.}
        Im Schnitt 34\,\% der Laufzeit blockiert, schlimmstenfalls 72\,\%.
        In NebulaStream würde das die gesamte Pipeline aufstauen.

  \item \textbf{Cache-Trefferquote skaliert linear.}
        8\,MB $\rightarrow$ 5\,\%, 32\,MB $\rightarrow$ 19\,\%, 128\,MB $\rightarrow$ 64\,\%.
        32\,MB ist ein guter Kompromiss.
\end{enumerate}

\subsection*{Was wir daraus mitnehmen}

\begin{itemize}[nosep]
  \item RocksDB \emph{kann} auf 256\,MB laufen -- aber nur mit konservativer
        Konfiguration und aktiver Kompression.
  \item Die Konfigurationswahl ist wichtiger als die Hardware. NebulaStream
        braucht durchdachte Defaults oder Auto-Tuning.
  \item Write-Stalls sind ein reales Risiko für die Streamverarbeitung und
        müssen im Pipeline-Design berücksichtigt werden.
\end{itemize}

\subsection*{Einschränkung}

Nur ein Workload, drei von fünf Parametern, synthetische Datengrößen,
simulierte Hardware. Die \emph{Richtung} der Ergebnisse ist belastbar, die
exakten Zahlen nicht direkt auf NebulaStream übertragbar.

\subsection*{Nächste Schritte}

\begin{enumerate}[nosep]
  \item Reinen Schreib-Workload (\texttt{fillrandom}) ergänzen
  \item Vollständigen Sweep (900 Kombinationen) mit allen 5 Parametern
  \item NebulaStream-realistische Schlüssel-/Wertgrößen testen
\end{enumerate}

\end{document}
